\section{Banach空间}
\label{sec:nspace-banach-spaces}

迭代算法、级数截断和逐步加细的逼近，常常只告诉我们近似之间越来越接近，即得到一个柯西列，而不直接给出极限。要从“彼此接近”推出“存在极限”，空间里必须真有这个极限。范数规定了怎样衡量接近，却不保证这一点；后者是空间与范数共同的性质。

\begin{definition}[Banach空间]
\label{def:nspace-banach}
若赋范空间$V$中的每个柯西列都依范数收敛到$V$中的某个元素，即范数诱导的度量空间是完备的，称$V$是\term{完备}（complete）的赋范空间，也称\term{Banach空间}（Banach space）。
\end{definition}

这就是第~\ref{chap:metric-spaces}章的完备性，只是度量来自范数。卷二已经证明有限维赋范空间在任意范数下都完备，所以有限维时无需另作核验。无限维时，完备与否要逐个空间证明；同一个集合换一个范数，结论可能改变，第~\ref{sec:nspace-continuous-function-spaces}节会给出这样的例子。

\paragraph{闭子空间判据}
很多函数空间是另一个已知完备空间的子空间。这时可以借用第~\ref{chap:metric-spaces}章的两个结论：完备度量空间中的闭子集是完备的；任何度量空间中的完备子集都是闭集。用在Banach空间$X$的线性子空间$M$上，就得到下面的判据。

\begin{theorem}[闭子空间判据]
\label{thm:nspace-closed-subspace}
设$X$是Banach空间，$M\subseteq X$是线性子空间，$M$带有限制后的范数。则$M$是Banach空间当且仅当$M$是$X$中的闭集。
\end{theorem}

\begin{proof}
设$M$闭，$(f_n)$是$M$中的柯西列。它也是$X$中的柯西列，由$X$完备收敛到某个$f\in X$；$M$闭，所以$f\in M$。反过来，设$M$完备，$f_n\in M$且在$X$中$f_n\to f$。收敛列是柯西列，所以$(f_n)$是$M$中的柯西列，在$M$中收敛到某个$g$；$X$中极限唯一，$f=g\in M$，所以$M$闭。
\end{proof}

两个方向的分工不同。从闭推出完备时，大空间$X$的完备性提供极限，$M$的闭性保证极限没有离开$M$。从完备推出闭时，并不需要$X$完备，第~\ref{chap:functions-as-vectors}章证明有限维子空间闭性时用的正是这一方向。

\paragraph{不闭的子空间}
多项式给出一个不完备的例子。在$[0,1]$上用$\lVert f\rVert_\infty=\max_x\lvert f(x)\rvert$衡量多项式，取$e^x$的Taylor多项式$p_N(x)=\sum_{k=0}^Nx^k/k!$。由Lagrange余项，
\[
 \max_{x\in[0,1]}\lvert e^x-p_N(x)\rvert\leq\frac{e}{(N+1)!}\to0,
\]
所以对$m>n$，$\lVert p_m-p_n\rVert_\infty\leq\frac{e}{(m+1)!}+\frac{e}{(n+1)!}$，$(p_N)$是多项式空间中的柯西列。

它在多项式空间中没有极限。若有多项式$p$使$\lVert p_N-p\rVert_\infty\to0$，则对每个$x$都有$p_N(x)\to p(x)$；又$p_N(x)\to e^x$，所以$p=e^x$在$[0,1]$上成立。但多项式求足够多次导数后恒为零，$e^x$的各阶导数都不为零，矛盾。所以多项式空间在最大误差下不完备；下一节证明$C[0,1]$完备后，闭子空间判据说明这等价于多项式空间在$C[0,1]$中不闭。

\paragraph{范数与完备性是两个问题}
给一个函数集合配范数，回答的是“怎样衡量误差”；问它是否完备，回答的是“误差趋零的过程能否在集合内找到终点”。前者只要核验三条公理，后者要分析柯西列的极限。不完备的赋范空间仍然是合法的赋范空间，只是不能仅凭柯西性断言极限存在。

不完备的空间可以通过补入极限变成完备的：第~\ref{chap:metric-spaces}章的完备化构造对赋范空间同样适用，线性运算和范数都可以沿柯西列延拓。但抽象完备化只保证新元素存在，不说明它们是哪一类函数。本章采用另一条路：直接在更大的函数集合上定义范数，再证明这个具体的空间完备，第~\ref{sec:nspace-lp-spaces}节的$L^p$就是这样得到的。
